% SPDX-License-Identifier: LPPL-1.3c
% Copyright (C) 2026 Christophe Jorssen
% Author and Current Maintainer: Christophe Jorssen <christophe.jorssen@gmail.com>
% LPPL maintenance status: maintained
% This file is part of luafemm. See LICENSE and LICENSES.md for its terms.
\section{Extended electrostatic models and measurements}
\label{sec:electric-extended}

These facilities extend |model=finite elements| in both planar and
axisymmetric geometry. They also apply to |model=source correction| within
the source guards and interaction-energy conventions of
Section~\ref{sec:electric-hybrid}. Ordinary field-energy quantities do not
extend to ideal singular sources. Declare the domain, dielectric regions, electrodes
and exterior conditions first. Rendering a field, capturing a profile,
measuring a force or calculating a capacitance matrix solves and freezes the
physical model. A mesh pic already freezes its geometry. Afterwards only
views, measurements and drawing styles may change. A named matrix or force
owns its data: reusing a picture name does not alter a previous measurement.

\subsection{Equipotential conductors with prescribed charge}
\begin{key}{/femm/electric electrode/condition=\meta{voltage or charge} (initially voltage)}
Path-local physical choice. |voltage| requires |potential| in volts.
|charge| introduces one unknown equipotential and imposes its total free
charge. Do not also supply |potential|. At least one prescribed potential or
a positive mixed boundary coefficient must supply a reference, unless an
antiperiodic constraint removes the constant mode. A neutral conductor alone
does not supply an absolute reference.
\end{key}
\begin{key}{/femm/electric electrode/charge=\meta{number} (initially empty)}
Total signed charge in coulombs, only for |condition=charge|. Empty means zero:
a neutral floating conductor. In a planar model this is the charge over the
specified |depth|, not C/m. In axisymmetry it is the charge of the complete
revolved body. For a prescribed voltage, even an explicit zero charge is
rejected as a second physical condition.
\end{key}
\begin{key}{/femm/electric electrode/name=\meta{text} (initially empty)}
An empty name receives an automatically generated name. Several paths with the
same explicit name form one conductor, including separated pieces connected
by an ideal wire outside the section. Their condition and prescribed value
must agree. The charge applies once to the whole named conductor; it is not
repeated for each piece. The |interior| setting remains local to each path.
Distinct named electrodes may not touch. A charge-controlled electrode may
not touch a voltage boundary, including the default grounded box; use an
explicit insulating side where appropriate. |outer boundary| is reserved.
\end{key}
The usual electrode |voltage| result now returns the solved potential of a
charge-controlled conductor. Its |charge| result is the variational reaction,
which can be compared with the prescribed total charge. Charges are obtained
from the assembled equations, not from a smoothed boundary field.

\subsection{A Maxwell capacitance matrix}
\begin{command}{\femmelectricmatrix\oarg{options}}
Capture a matrix after all physical declarations.
Its key family is |/femm/electric matrix|. For the selected voltage electrodes,
\[
 Q_i=q_i^{(0)}+\sum_j C_{ij}V_j.
\]
The offset includes volume/surface sources, specified floating charges and
unselected electrode voltages. Each matrix column uses a unit voltage and
zero independent sources on the same owned mesh. Floating conductors are
eliminated with zero incremental charge. The result is a Maxwell matrix:
off-diagonal entries can be negative. It is not a table of positive two-terminal
capacitances, nor does its computation add a circuit equation to the model.
\end{command}
\begin{key}{/femm/electric matrix/name=\meta{text} (initially empty)}
Required nonempty document-lifetime result name. Successful reuse replaces
that named matrix atomically. A failed computation does not replace it.
\end{key}
\begin{key}{/femm/electric matrix/model=\meta{snapshot} (initially empty)}
Empty selects the current electric picture. Otherwise use a previously named
electric snapshot, including outside a picture. The source state is unchanged.
\end{key}
\begin{key}{/femm/electric matrix/electrodes=\meta{comma-separated names} (initially empty)}
Empty selects all voltage-controlled named electrodes in declaration order.
An explicit list fixes the row/column order. Quote the list with braces.
Unknown, duplicate or charge-controlled ports are errors. Unselected voltage
electrodes remain fixed in the offset and are grounded in unit-column solves.
\end{key}
\begin{key}{/femm/electric matrix/normalization=\meta{total or per length} (initially total)}
|total| gives coefficients in farads and offsets in coulombs. |per length|
gives F/m and C/m for a planar model. It is rejected in axisymmetry. Reciprocity
is checked in the numerical validation; discretization and algebraic tolerance
still limit accuracy. A matrix can be singular when its electrodes share a
free common-mode voltage, even though each unit experiment has a reference.
\end{key}
\begin{command}{\femmelectricmatrixvalue\oarg{options}}
Typeset one matrix entry or source offset with the following keys.
\end{command}
\begin{key}{/femm/electric matrix value/matrix=\meta{name} (initially empty)}
Required captured matrix name; no current-picture dependency.
\end{key}
\begin{key}{/femm/electric matrix value/row=\meta{electrode} (initially empty)}
Required receiving electrode name, selected in the captured matrix.
\end{key}
\begin{key}{/femm/electric matrix value/column=\meta{electrode} (initially empty)}
Excited electrode name; required for a capacitance entry, ignored for an offset.
\end{key}
\begin{key}{/femm/electric matrix value/quantity=\meta{capacitance or offset charge} (initially capacitance)}
Choose $C_{ij}$ or $q_i^{(0)}$. Units follow the captured normalization.
\end{key}
\begin{key}{/femm/electric matrix value/number format=\meta{PGF number options} (initially sci,precision=3)}
Formatting only, using native |\pgfmathprintnumber| options.
\end{key}
\begin{command}{\femmelectricmatrixplot\oarg{options}}
Inside a PGFPlots axis, plot a matrix column against integer row numbers
$1,\ldots,n$. Supply electrode tick labels and units through ordinary axis
keys. The coefficient values retain the captured SI normalization.
\end{command}
\begin{key}{/femm/electric matrix plot/matrix=\meta{name} (initially empty)}
Required matrix name. Each plot reads the owned result without another solve.
\end{key}
\begin{key}{/femm/electric matrix plot/column=\meta{electrode} (initially empty)}
Required column electrode name.
\end{key}
\begin{key}{/femm/electric matrix plot/plot options=\meta{PGFPlots options} (initially empty)}
Passed to |\addplot+|. Use normal axis styles for bars, colours, markers and legends.
The worked example in Section~\ref{sec:tutorial-electric-floating} shows all
three commands with a neutral intermediate screen.
\end{key}

\subsection{Positive symmetric dielectric tensors}
In a finite-element calculation, the constitutive law is
$\mathbf D=\epsilon_0\boldsymbol\epsilon_r\mathbf E$ with
\[
 \boldsymbol\epsilon_r=R(\theta)
 \begin{pmatrix}\epsilon_x&0\\0&\epsilon_y\end{pmatrix}R(\theta)^T,
 \qquad \epsilon_x>0,\quad\epsilon_y>0.
\]
The solver uses this tensor in stiffness, displacement and
$\tfrac12\mathbf E\mathbin\cdot\mathbf D$ energy. Field-line pics follow
$\mathbf E$, which need not be parallel to $\mathbf D$.
\begin{key}{/femm/electrostatic/relative permittivity x=\meta{positive number} (initially empty)}
Picture-local first principal relative permittivity, dimensionless.
Empty uses the scalar |relative permittivity| or selected material value.
\end{key}
\begin{key}{/femm/electrostatic/relative permittivity y=\meta{positive number} (initially empty)}
Second principal value with the same default rule. These are principal values
before rotation, not the diagonal entries after rotation.
\end{key}
\begin{key}{/femm/electrostatic/material angle=\meta{degrees} (initially empty)}
Counterclockwise orientation of the first principal axis relative to physical
$x$; empty means zero. TikZ drawing transformations do not rotate this
constitutive input. Analytical models still require an isotropic medium.
\end{key}
\begin{key}{/femm/electric region/relative permittivity x=\meta{positive number} (initially empty)}
Path-local first principal value. A region with no constitutive input inherits
the complete background tensor. When specifying principal values or an angle,
missing principal values inherit the background principal values.
\end{key}
\begin{key}{/femm/electric region/relative permittivity y=\meta{positive number} (initially empty)}
Path-local second principal value, with the same inheritance rule.
\end{key}
\begin{key}{/femm/electric region/material angle=\meta{degrees} (initially empty)}
Path-local orientation. A constitutive override starts at angle zero unless
an explicit angle follows. A scalar |relative permittivity| assignment resets
both principal overrides and the angle. A catalogue selection resets both
principal values and sets angle zero; subsequent keys override them in order.
In axisymmetry $(x,y)=(r,z)$. A tensor touching the regular axis must have zero
$r$--$z$ coupling; an incompatible rotated tensor is rejected. This is a
meridional tensor model, not general three-dimensional anisotropy.
\end{key}

\begin{codeexample}[width=0pt]
\begin{tikzpicture}[femm/electrostatic={model=finite elements,
    xmin=0,xmax=2,ymin=0,ymax=1,outer boundary=insulated,
    relative permittivity x=2,relative permittivity y=4,
    material angle=45,mesh size=.15},scale=1.8]
  \draw[thick,femm/electric electrode={potential=10}] (0,0)--(0,1);
  \draw[thick,femm/electric electrode={potential=0}] (2,0)--(2,1);
  \pic {femm electric field};\pic {femm equipotentials};
  \node[below] at (1,0) {$\theta=45^\circ$};
\end{tikzpicture}
\end{codeexample}
The insulating horizontal sides impose $D_y=0$, not $E_y=0$. Thus this example
has a genuinely two-dimensional field near the electrode corners.

\subsection{Mixed and translated periodic boundaries}
\begin{key}{/femm/electric boundary/condition=\meta{insulated, voltage, displacement, mixed, periodic or antiperiodic} (initially insulated)}
The additional choices apply to physical paths on the rectangular exterior.
For |mixed| the convention is $\mathbf D\cdot\mathbf n=cV+g$, with outward
normal. This adds $\int cN_iN_j$ to stiffness and $-\int gN_i$ to the load.
The regular axis continues to accept only |insulated|. Mixed or periodic paths
may not overlap an electrode edge.
\end{key}
\begin{key}{/femm/electric boundary/coefficient=\meta{nonnegative number} (initially 0)}
$c$ in F/m$^2$, for |condition=mixed|. A positive value supplies a potential
reference; zero reduces to prescribed displacement. Negative coefficients
are rejected to preserve the positive variational problem.
\end{key}
\begin{key}{/femm/electric boundary/normal displacement=\meta{number} (initially empty)}
For |displacement| this is the required outward $D_n$ in C/m$^2$.
For |mixed| it is $g$ in C/m$^2$, default zero. A nonzero mixed term therefore
need not mean a nonzero outward flux: the $cV$ term must also be included.
\end{key}
\begin{key}{/femm/electric boundary/pair=\meta{text} (initially empty)}
Required for |periodic| or |antiperiodic|. Exactly two straight, open paths
must share the name and condition, lie on opposite box sides, and have
identical oriented vectors. Both paths therefore run in the same drawing
direction. They impose respectively $V_2=V_1$ or $V_2=-V_1$; the natural
normal displacement has the corresponding opposite outward sign.

Boundary subdivisions are synchronized before assembly. Incompatible corner
voltages and incomplete pairs fail explicitly. Very aggressive refinement
which does not synchronize within twelve mesh passes also fails; simplify the
boundary refinement in that case. A purely periodic insulated cell still
needs a reference. Antiperiodicity removes the constant mode. In axisymmetry
only axial translations are accepted, since radial translations change the
integration measure. Rotational, curved and sliding periodic interfaces are
not part of this electric boundary contract.
\end{key}
\begin{key}{/femm/electric value/quantity=\meta{boundary energy, total energy or periodic flux}}
Additional global quantities, alongside the existing choices (whose default
remains |charge|). |boundary energy| is $\tfrac12\int cV^2\,\mathrm dS$ in J;
|total energy| adds it to the dielectric field energy. |periodic flux| is the
outward displacement flux in C represented by reactions on paired nodes
without an electrode owner. The global |charge balance| includes this flux;
antiperiodic cells need not have zero ordinary outward flux. All three results
include planar depth or the full revolution. They are not field-energy
per-length aliases.
\end{key}

\Needspace{11cm}
\begin{codeexample}[width=0pt]
\begin{tikzpicture}[femm/electrostatic={model=finite elements,
    xmin=0,xmax=2,ymin=0,ymax=1,outer boundary=insulated,
    mesh size=.2},scale=1.8]
  \draw[thick,femm/electric electrode={potential=10}] (0,0)--(0,1);
  \draw[dashed,femm/electric boundary={condition=mixed,
      coefficient=8.8541878128e-10}] (2,0)--(2,1);
  \path[femm/electric boundary={condition=periodic,pair=cell}]
    (0,0)--(2,0);
  \path[femm/electric boundary={condition=periodic,pair=cell}]
    (0,1)--(2,1);
  \pic {femm electric field};\pic {femm equipotentials};
\end{tikzpicture}
\end{codeexample}
The right side is a capacitive Robin termination, not a grounded electrode.
The top and bottom represent translated copies of the same cell.

\subsection{Forces and torque from a Maxwell contour}
\begin{key}{/tikz/femm/electric force=\meta{options}}
Measurement on one simple closed path, after all physical declarations.
The contour must stay in one homogeneous, isotropic, source-free dielectric,
and must not touch or cross an electrode, charged sheet or material interface.
Its interior may enclose those bodies. Both orientations give the outward
traction. Ordinary TikZ drawing options control whether/how the contour is shown.
The force is $\oint[\epsilon\mathbf E\mathbf E^T-
\tfrac12\epsilon E^2I]\mathbf n\,\mathrm dS$.
\end{key}
\begin{key}{/femm/electric force/name=\meta{text} (initially empty)}
Required document-lifetime result name; successful reuse replaces the result.
\end{key}
\begin{key}{/femm/electric force/origin x=\meta{number} (initially 0)}
Planar torque origin in physical model coordinates and the declared length unit.
This does not translate the contour or change its force.
\end{key}
\begin{key}{/femm/electric force/origin y=\meta{number} (initially 0)}
Second torque-origin coordinate. Origins are irrelevant to the axisymmetric
axial-force result; no meridional ``radial resultant'' is offered as a 3D force.
\end{key}
\begin{key}{/femm/electric force/curve tolerance=\meta{positive number} (initially .001)}
Flattening tolerance in model units, distinct from mesh size. Integration splits
each straight piece at triangle faces and uses raw element fields. An axis
closure has zero revolved area and contributes no traction. The model's
|boundary snap tolerance| also preserves this closure after TikZ coordinate
rounding and transformations.
\end{key}
\begin{command}{\femmelectricforcevalue\oarg{options}}
Typeset a previously captured force or torque; no solve or current picture is needed.
\end{command}
\begin{key}{/femm/electric force value/force=\meta{name} (initially empty)}
Required captured force name.
\end{key}
\begin{key}{/femm/electric force value/quantity=\meta{Fx, Fy, Fz, torque, Fx per length, Fy per length or torque per length} (initially Fy)}
Planar total |Fx|, |Fy| in N and |torque| in N\,m include depth; their per-length
variants use N/m and N. Positive torque is counterclockwise about the specified
origin. Axisymmetry offers only |Fz| in N for the full revolution. Unavailable
quantities are rejected; explicitly select |Fz| for an axisymmetric result.
\end{key}
\begin{key}{/femm/electric force value/number format=\meta{PGF number options} (initially sci,precision=3)}
Native number formatting, independent of physical accuracy.
\end{key}

\Needspace{13cm}
\begin{codeexample}[width=0pt]
\begin{tikzpicture}[femm/electrostatic={model=finite elements,
    xmin=0,xmax=2,ymin=0,ymax=1,outer boundary=insulated,
    mesh size=.05,name=force cell},scale=2.2]
  \draw[thick,femm/electric electrode={potential=100}] (0,0)--(0,1);
  \draw[thick,femm/electric electrode={potential=0}] (2,0)--(2,1);
  \draw[red,very thick,femm/electric charge layer={
      surface charge density=2e-8}] (1,.3)--(1,.7);
  \pic {femm electric field};
  \draw[green!40!black,dashed,femm/electric force={name=sheet force}]
    (.6,.1) rectangle (1.4,.9);
\end{tikzpicture}
\qquad
$F_x=\femmelectricforcevalue[force=sheet force,quantity=Fx]$ N.
\end{codeexample}
With the default centimetre unit and one-centimetre depth, the imposed uniform
bias is $5000$ V/m and the sheet charge is $8\times10^{-13}$ C. Symmetry cancels
the sheet's net self-force, giving $F_x=4\times10^{-9}$ N in the continuum.
Compare refinements and move the contour away from the sheet ends: raw P1
stress integrals need not converge monotonically on a fixed contour. A small
linear residual or charge-balance error is not a force-accuracy certificate.
This method does not model electrostriction or stress inside an anisotropic body.

\subsection{Native FEMM point properties versus Coulomb nodes}
The FEMM~4 and FEMM~5 electrostatic point-property parameter |qp| is a line
charge density in C/m. In planar geometry it represents a line perpendicular
to the section; in axisymmetry a node at radius $r>0$ represents a ring with
total charge $2\pi r\,q_p$. In the audited native implementation the radial
weight vanishes at $r=0$; it does not create an isolated Coulomb charge there.
A zero |qp| selects prescribed point potential rather than a removable zero
source. Such nodal point constraints and singular loads require their own
mesh-limit interpretation.

Consequently |femm/point charge| remains a three-dimensional Coulomb source
for the analytical provider, and is rejected by the FE provider. No silent
conversion is made. To describe a finite charged body use volume charge or a
charged interface; to describe a finite conductor use an electrode. Native
|.fee|/|.res| interchange and native point properties are handled by the separate
file adapter in Section~\ref{sec:electric-exchange}.
